\section{从扩散模型到流模型：范式转换}

\subsection{时间方向的翻转}

在正式引入 Flow Matching 之前，讲者特别强调了一个\textbf{极易混淆}的符号约定变化：

\begin{warningbox}{时间方向的翻转——最容易混淆的地方}
\textbf{扩散模型}的时间方向（前向过程视角）：
\begin{itemize}
    \item $t = 0$：数据分布 $p_{\text{data}}$
    \item $t = T$：噪声分布（标准高斯）$\mathcal{N}(0, I)$
\end{itemize}

\textbf{Flow Matching}的时间方向（生成过程视角）：
\begin{itemize}
    \item $t = 0$：噪声分布（标准高斯）$p_0 = \mathcal{N}(0, I)$
    \item $t = 1$：数据分布 $p_1 = p_{\text{data}}$
\end{itemize}

也就是说，$x_0$ 在 Flow Matching 中代表高斯噪声样本，$x_1$ 代表真实数据样本。这与扩散模型的习惯完全相反！
\end{warningbox}

讲者坦言，即使是他本人也经常被这两套符号搞混。之所以保留这种不一致，是因为扩散模型和流匹配分别来自不同的文献传统，各自的论文和教材都遵循各自的约定。

\subsection{扩散模型与流模型的关键区别}

\begin{importantbox}{核心区别：参考分布的灵活性}
\textbf{扩散模型}中，所有推导都假设参考分布（噪声分布）是标准高斯分布 $\mathcal{N}(0, I)$。如果更换为其他分布（如各向异性高斯），所有数学推导都需要重新进行。

\textbf{流模型}的一大优势是：参考分布 $p_0$ 可以是\textbf{任意分布}。无论 $p_0$ 怎么变化，流匹配的数学公式保持不变。这使得流模型可以被看作扩散模型的一种\textbf{泛化}。
\end{importantbox}

尽管两者来自不同的文献传统，最终它们在数学上是等价的（在标准高斯参考分布的特殊情况下）。这也是本讲要揭示的一个重要洞见。

\subsection{本章小结}

Flow Matching 从生成过程的视角出发，将时间方向翻转，并允许任意参考分布。虽然最终可以与扩散模型建立等价关系，但这种新视角带来了更简洁的训练目标和更灵活的模型设计空间。

